% META: {"id": "TSPE-DERCONV-CR-014", "chapitre": "TSPE-DERIVATION-CONVEXITE", "type_objet": "cours", "capacites_codes": ["C5"], "sources_inspiration": [], "status": "approved"}

\section{Lecture graphique de la convexité}

% ============================================================
% STRATE 1 — Methodes de lecture
% ============================================================

\propriete{
A partir de la représentation graphique de $f$ :
\begin{itemize}
  \item $f$ est \textbf{convexe} sur un intervalle $I$ si la courbe est au-dessous de chaque corde reliant deux points de la courbe sur $I$.
  \item $f$ est \textbf{concave} sur un intervalle $I$ si la courbe est au-dessus de chaque corde.
  \item Un \textbf{point d'inflexion} est un point ou la courbe change de position par rapport a ses tangentes (elle passe « d'un cote a l'autre »).
\end{itemize}
}

\propriete{
A partir de la représentation graphique de $f'$ :
\begin{itemize}
  \item $f$ est convexe sur $I$ si et seulement si $f'$ est \textbf{croissante} sur $I$.
  \item $f$ est concave sur $I$ si et seulement si $f'$ est \textbf{décroissante} sur $I$.
  \item Les points d'inflexion de $f$ correspondent aux extremums de $f'$.
\end{itemize}
}

\propriete{
A partir de la représentation graphique de $f''$ :
\begin{itemize}
  \item $f$ est convexe sur les intervalles ou $f'' \geq 0$ (courbe de $f''$ au-dessus de l'axe des abscisses).
  \item $f$ est concave sur les intervalles ou $f'' \leq 0$.
  \item Les points d'inflexion de $f$ correspondent aux changements de signe de $f''$.
\end{itemize}
}

\exemple{
On observe le graphique de $f'$ : $f'$ est croissante sur $]-\infty, 2[$ et décroissante sur $]2, +\infty[$.

On en deduit que $f$ est convexe sur $]-\infty, 2[$ et concave sur $]2, +\infty[$, avec un point d'inflexion d'abscisse $x = 2$.
}

% ============================================================
% STRATE 2 — Applications
% ============================================================

\margeAppui{
\textbf{Correspondances à retenir :} $f''>0$ : $f'$ croissante, $f$ convexe. $f''<0$ : $f'$ décroissante, $f$ concave. $f''=0$ avec changement de signe : extremum de $f'$, point d'inflexion de $f$.
}

\erreurFrequente{
\textbf{Confondre le graphe de $f$ et celui de $f'$.} Sur le graphe de $f'$, les zeros donnent les extremums de $f$, et les extremums de $f'$ donnent les points d'inflexion de $f$.
}

% BEGIN-VERIFY
% from sympy import symbols, diff, Rational
%
% x = symbols('x', real=True)
%
% # Convexe <=> f' croissante. Sur x^2 :
% f = x**2
% fp = diff(f, x)
% echantillon = [Rational(k, 2) for k in range(-8, 9)]
% valeurs = [fp.subs(x, t) for t in echantillon]
% assert valeurs == sorted(valeurs)
% # Et la corde est bien au-dessus de la courbe.
% g = lambda t: t**2
% for k in range(0, 11):
%     t = Rational(k, 10)
%     assert g(t*(-2) + (1 - t)*3) <= t*g(-2) + (1 - t)*g(3)
%
% # Concave <=> f' decroissante. Sur -x^2 :
% h = -x**2
% hp = diff(h, x)
% valeurs_h = [hp.subs(x, t) for t in echantillon]
% assert valeurs_h == sorted(valeurs_h, reverse=True)
%
% # Les points d'inflexion de f correspondent aux extremums de f'. Sur x^3 :
% k = x**3
% kp = diff(k, x)                      # 3x^2, minimum en 0
% assert diff(kp, x).subs(x, 0) == 0
% assert kp.subs(x, 0) == 0
% assert all(kp.subs(x, t) > 0 for t in (-1, Rational(-1, 2), Rational(1, 2), 1))
% # f'' change bien de signe en 0 : c'est un point d'inflexion.
% assert diff(k, x, 2).subs(x, -1) < 0 < diff(k, x, 2).subs(x, 1)
% END-VERIFY
